\section{Illustration: the Philippine coupled run}
\label{sec:results}

We illustrate the implemented forward pass on the Philippine calibration (\OGCore{} via the
\texttt{ogphl} multi-sector build, $\Mind=8$ industries, $\Igood=5$ goods; electricity is
industry $m_e=2$), coupled to a \CLEWS{}/\OSeMOSYS{} scenario pair (a baseline and a
policy/decarbonisation reform). The application tests whether the principal forward
mappings are accepted by the economy solve and whether their responses are consistent
with the mechanisms of \Cref{sec:channels}; it is not a policy evaluation. Unless noted
otherwise, the aggregate results come from the committed M=8 golden record
(\texttt{results/golden.json}). Rates, wages, and quantities are reported as percent
changes from baseline. The sectoral capital decomposition is explicitly identified as an
earlier solve.

\subsection{The flagship coupled run}
\label{sec:coupled}
The \chan{coupled} experiment applies, in one reform: the \CLEWS{} electricity-cost
proxy (the ``auto'' source---a cost index or reconstructed levelized cost,
\Cref{sec:signal}) through the composite transmission of \Cref{sec:routes}
(inter-industry cost-push plus a recycled household wedge); public grid investment; and
the Global Burden of Disease health channel; then it emits, alongside the equilibrium
return and activity, a \$50/tCO\textsubscript{2} carbon-price input for the energy
model. The emitted inputs are inert in a single pass, so no carbon-price effect enters
the numbers reported here---the underlying \CLEWS{} reform is a policy-composition
scenario whose solve carries no emissions penalty, and the macro response reflects its
costs, investment, and emissions alone. The Philippine electricity-cost proxy rises about $5\%$ on average
(\Cref{sec:calibration}); the coupled run moves electricity output by $-2.0\%$.

The response is a positive transition and a negative steady state
(\Cref{tab:phl-coupled}): output is $+0.18\%$ near-term and $+0.20\%$ by period ten, as
the working-age morbidity gain raises effective labour, but $-0.14\%$ in the long-run
steady state, where the added retirees (the cleaner-air survival gain skews elderly)
combine with the contractionary cost channels (\Cref{fig:transition}). The isolated
health experiment supports the demographic interpretation of the horizon split; a model
without age structure cannot represent that cohort-composition mechanism directly.

\begin{table}[t]
\centering\footnotesize
\caption{The Philippine coupled run (percent change from baseline; \OGCore{} steady state
and transition periods $t_0$, $t_{10}$). Source: \texttt{results/golden.json}, run
\chan{coupled}.}
\label{tab:phl-coupled}
\begin{tabular}{@{}lrrr@{}}
\toprule
 & steady state & $t_0$ (near term) & $t_{10}$ \\
\midrule
output $Y$        & $-0.138$ & $+0.179$ & $+0.202$ \\
consumption $C$   & $-0.087$ & $+0.066$ & $+0.385$ \\
capital $K$       & $-0.128$ & $+0.164$ & $+0.218$ \\
labour $L$        & $-0.007$ & $+0.276$ & $+0.222$ \\
return $r$        & $-0.032$ & $+0.047$ & $-0.056$ \\
wage $w$          & $-0.117$ & $-0.117$ & $+0.009$ \\
\bottomrule
\end{tabular}
\end{table}

\begin{figure}[t]
\centering
\includegraphics[width=0.82\linewidth]{coupled_transition}
\caption{The coupled transition path (percent change from baseline): the near-term path
is positive, driven by the working-age morbidity gain, while the steady-state effect is
negative ($-0.14\%$). Source: \texttt{results/figures/coupled\_transition.png}.}
\label{fig:transition}
\end{figure}

\subsection{Standalone mechanism experiments}
\label{sec:decomposition}
Table~\ref{tab:phl-channels} reports \emph{runnable experiments}, each isolating one
mechanism as a standalone reform---so \chan{energy\_price} is the household wedge of
the energy channel alone, using the cost proxy (the composite is what \chan{coupled}
transmits), and \chan{clean\_incidence} is its recycled variant;
\Cref{fig:channel-impact} plots the steady-state output impacts. Four things stand out,
all as \Cref{sec:channels} anticipates.
The household-wedge \chan{energy\_price} experiment is small ($Y_{ss}=-0.006\%$): the
Philippine electricity-cost proxy rises about $5\%$ and produces little aggregate-output
movement through this route. The larger controlled $+20\%$ shocks are transmission
stress tests rather than application magnitudes. \emph{Carbon} is the clearest standalone contraction ($Y_{ss}=-0.021\%$): a
resource cost that shrinks capital and labour. \emph{Capital intensity} barely moves
aggregate output ($+0.001\%$) but reshapes the energy sector (\Cref{sec:ch-capint}).
\emph{Health} is near-zero at the steady state ($-0.0003\%$) but $+0.19\%$
near-term---the morbidity productivity gain.
The forward emitters (\chan{discount\_rate}, \chan{demand}) and \chan{investment} are
$\approx 0$ as standalone reforms: the Philippine public grid-capex delta is small, and
the emitters only feed \CLEWS{} until the loop closes. Finally, the coupled steady state
($-0.138\%$) is not reproduced by summing these isolating experiments
($\approx-0.02\%$)---nor should it be: the isolating configurations do not collectively
duplicate the application treatment (the energy price, for one, enters the standalone
row as the household wedge alone but enters \chan{coupled} as the composite). General-
equilibrium interaction is possible, but this comparison does not identify it. A matched cumulative
decomposition---a sequence of runs whose final step reproduces \chan{coupled}---is the
instrument for attributing the difference; we do not attribute it here.

\begin{table}[t]
\centering\footnotesize
\caption{Standalone mechanism experiments: steady-state response (percent change from
baseline), one isolating reform per row, on the M=8 Philippine calibration.
Source: \texttt{results/golden.json}.}
\label{tab:phl-channels}
\begin{tabular}{@{}lrrrrr@{}}
\toprule
experiment & $Y$ & $C$ & $K$ & $L$ & $r$ \\
\midrule
\chan{energy\_price} (cost proxy)      & $-0.006$ & $-0.062$ & $-0.011$ & $+0.009$ & $+0.002$ \\
\chan{clean\_incidence} (recycled)    & $+0.003$ & $-0.092$ & $-0.003$ & $+0.017$ & $+0.017$ \\
\chan{carbon}                         & $-0.021$ & $+0.026$ & $-0.019$ & $-0.045$ & $-0.064$ \\
\chan{capital\_intensity}             & $+0.001$ & $+0.144$ & $+0.012$ & $-0.018$ & $-0.013$ \\
\chan{investment}                     & $+0.000$ & $+0.000$ & $+0.000$ & $+0.000$ & $-0.000$ \\
\chan{health}                         & $-0.000$ & $+0.024$ & $+0.004$ & $-0.013$ & $-0.011$ \\
\midrule
\chan{coupled} (application config.)  & $-0.138$ & $-0.087$ & $-0.128$ & $-0.007$ & $-0.032$ \\
\bottomrule
\end{tabular}
\end{table}

\begin{figure}[t]
\centering
\includegraphics[width=0.82\linewidth]{golden_channel_impact}
\caption{Steady-state output impact of the standalone mechanism experiments on the M=8
Philippine calibration (the figure's ``by channel'' title predates the
channel/experiment distinction of \Cref{sec:decomposition}).
Source: \texttt{results/figures/golden\_channel\_impact.png}.}
\label{fig:channel-impact}
\end{figure}

\subsection{Distributional incidence}
\label{sec:incidence}
Under the recycled
electricity-cost change with energy a partial necessity (\chan{clean\_incidence}),
aggregate consumption falls ($C_{ss}=-0.092\%$) even as the price change is returned as
transfers: the burden lands more heavily on lower-income households, who spend a larger
share on energy---a regressive pattern a single-agent macro core cannot produce.

\subsection{The capital pair}
\label{sec:capital-results}
The \chan{capital\_intensity} channel is decomposed exactly via the firm identity of
\Cref{sec:ch-capint} (\Cref{tab:capital-pair}, measured on an earlier steady-state solve
of the lever): raising the energy capital exponent
$12.2\%$ collapses the electricity price $42.8\%$; demand is elastic enough that output
\emph{rises} $20.9\%$, yet capital \emph{falls} $22.4\%$ because the price collapse shrinks
revenue---and the identity $\hat K_E=\hat\gamma_E\hat p_E\hat Y_E$ reproduces that response
to within $0.03\%$. It is a factor-share and energy-price lever, not a capital-draw-in.

\subsection{Validation}
\label{sec:validation}
Validation is layered because the tests establish different claims
(\Cref{tab:validation-layers}). The grouped run battery executes the link in its own
environment and the country \OGCore{} model in the country model's environment,
exchanging data files rather than sharing an interpreter. All battery items pass the
solve-completion gate; selected macro aggregates are regression-locked in
\texttt{results/golden.json} and re-checked automatically.

\begin{table}[t]
\centering\footnotesize
\caption{Validation layers and their current boundaries.}
\label{tab:validation-layers}
\begin{tabularx}{\linewidth}{@{}>{\bfseries\raggedright\arraybackslash}p{0.20\linewidth} L L@{}}
\toprule
Evidence & What it establishes & What it does not establish \\
\midrule
Transform and guard tests & dimensions, bounds, concordances, and failure behaviour & economic calibration or equilibrium response \\
Controlled mechanism runs & solve completion, response direction, and selected identities & joint attribution or application magnitude \\
Golden aggregate record & reproducibility of selected macroeconomic outputs & every sectoral or emitted interface artifact \\
Coupled forward run & joint solution of the price, public-investment, and health mappings & an additive decomposition or closed-loop equilibrium \\
\CLEWS{} re-solve test & a controlled demand patch survives copy, solve, and read-back & passage of emitted return and activity paths end to end \\
Closed-loop iteration & not yet tested & convergence and propagated-error control \\
\bottomrule
\end{tabularx}
\end{table}

One documented exception is that the lives-saved health reform leaves an intrinsic
Walras residual that scales
with the target, so its steady-state resource-constraint gate is loosened to $10^{-5}$
(an order of magnitude under \OGCore{}'s own transition-path default), while the
deaths-added direction stays at the tight default.

\caveat{The electricity-cost input is scenario-derived and ratio-based, and the health
burden and age profile are sourced from GBD; the translation from emissions to burden is
still calibrated. Carbon and investment \emph{levels} depend on the uncalibrated
unit/deflator bridge (\Cref{sec:calibration}), so treat those magnitudes as illustrative.}
